\documentclass[11pt,a4paper]{article}
\usepackage[T1]{fontenc}
\usepackage[utf8]{inputenc}
\usepackage{lmodern}
\usepackage[a4paper,margin=25mm]{geometry}
\usepackage{microtype}
\usepackage{amsmath,amssymb,mathtools,amsthm,bm}
\usepackage{booktabs}
\usepackage{enumitem}
\usepackage{algorithm}
\usepackage{algpseudocode}
\usepackage[numbers,sort&compress]{natbib}
\usepackage{xcolor}
\usepackage[colorlinks=true,allcolors=blue!55!black]{hyperref}
\usepackage[nameinlink,noabbrev]{cleveref}

\newtheorem{proposition}{Proposition}
\newtheorem{remark}{Remark}
\newcommand{\R}{\mathbb{R}}
\newcommand{\X}{\mathcal{X}}
\newcommand{\Y}{\mathcal{Y}}
\newcommand{\T}{\mathcal{T}}
\newcommand{\A}{\mathcal{A}}
\newcommand{\F}{\mathcal{F}}
\newcommand{\J}{\mathcal{J}}
\newcommand{\jump}[1]{\left[\!\left[#1\right]\!\right]}
\newcommand{\dd}{\,\mathrm{d}}
\newcommand{\abs}[1]{\left|#1\right|}

\title{Dual-Weighted-Residual Adaptive Loop\\
for Linear Time-Dependent PDEs}
\date{\today}


\begin{document}
\maketitle


% \section{Scope and Notation}

% The framework covers well-posed linear 1st-order-in-time equations that can be written in a Hilbert space as
% \begin{equation}
%   M\partial_t u+\mathcal{L}(t)u=f,\qquad t\in(0,T],
%   \label{eq:abstract-pde}
% \end{equation}
% with an initial condition and suitable boundary conditions.

% % Let
% % \[
% %   V\hookrightarrow H\simeq H^*\hookrightarrow V^*
% % \]
% % be a Gelfand triple. 
% Let \(m:H\times H\to\R\) be a continuous symmetric positive-definite mass form and let \(a(t;\cdot,\cdot):V\times V\to\R\) be a measurable continuous bilinear form satisfying the assumptions required for well-posedness. Define
% \begin{align*}
%   \X&:=\{v\in L^2(0,T;V):\partial_t v\in L^2(0,T;V^*)\},\\
%   \Y&:=\X.
% \end{align*}
% % The choice \(\Y=\X\) is slightly stronger than the minimal test space but keeps the trace \(v(0)\) used below well-defined. Elements of \(\X\) have continuous \(H\)-valued representatives and hence meaningful traces at \(t=0,T\) \citep{LionsMagenes1972}.



% \section{Continuous Primal and Adjoint Problems}

% For \(v\in\Y\), set
% \begin{align}
%   \A(u)(v)
%   &:={\int_0^T}\bigl[m(\partial_tu,v)+a(t;u,v)\bigr]\dd t
%        +m(u(0),v(0)),                                      \label{eq:A-cont}\\
%   \F(v)
%   &:={\int_0^T}\ell(t;v)\dd t+m(u_0,v(0)).                 \label{eq:F-cont}
% \end{align}
% The continuous primal problem is
% \begin{equation}
%   \boxed{\text{Find }u\in\X\text{ such that }\A(u)(v)=\F(v)
%   \quad\forall v\in\Y.}
%   \label{eq:continuous-primal}
% \end{equation}
% The initial condition is encoded weakly by the last terms in
% \cref{eq:A-cont,eq:F-cont}.

% Let the quantity of interest be $\J(u)$.
% % the continuous linear functional
% % \begin{equation}
% %   \J(u)=\int_0^T j(t;u(t))\dd t+j_T(u(T)).
% %   \label{eq:goal}
% % \end{equation}
% The continuous adjoint is
% \begin{equation}
%   \boxed{\text{Find }z\text{ such that }\A(w)(z)=\J(w)
%   \quad\forall w\in\X.}
%   \label{eq:continuous-adjoint}
% \end{equation}
% It is a terminal-value problem and is solved backward in time. This primal--adjoint construction and its goal-oriented identity are the foundation of the DWR method \citep{BeckerRannacher1996,BeckerRannacher2001,BangerthRannacher2003}.



% \section{Time and Space Discretisation}

% \subsection{Temporal Mesh and Jumps}

% Let
% \[
%   0=t_0<t_1<\cdots<t_N=T,\qquad I_n=(t_{n-1},t_n],
%   \qquad k_n=t_n-t_{n-1}.
% \]
% For a function that may be discontinuous at \(t_n\), write
% \[
%   v_n^-:=\lim_{t\uparrow t_n}v(t),\qquad
%   v_n^+:=\lim_{t\downarrow t_n}v(t),\qquad
%   [v]_n:=v_n^+-v_n^-.
% \]
% The discontinuous Galerkin time space of degree \(r\) is
% \begin{equation}
%   \X_k^r(V):=\{v:v|_{I_n}\in\mathbb{P}_r(I_n;V),\ 1\le n\le N\}.
%   \label{eq:time-space}
% \end{equation}
% This jump convention and the dG time formulation below follow the classical analysis of dG time stepping for parabolic equations \citep{ErikssonJohnsonThomee1985}.
% % check if this is ref [19] in 2021-..


% \subsection{Spatial Meshes}

% Let \(\T_h^n\) be a conforming spatial mesh on \(I_n\), and let \(V_h^n\subset V\) be its finite element space of degree \(p\). Define
% \begin{equation}
%   \X_{kh}^{r,p}:=
%   \{v:v|_{I_n}\in\mathbb{P}_r(I_n;V_h^n),\ 1\le n\le N\}.
%   \label{eq:full-space}
% \end{equation}
% The broken dG bilinear form is
% \begin{align}
%   \A_{kh}(U)(v)
%   :={}&\sum_{n=1}^N\int_{I_n}
%        \bigl[m(\partial_tU,v)+a(t;U,v)\bigr]\dd t \notag\\
%       &+\sum_{n=2}^N m([U]_{n-1},v_{n-1}^+)
%        +m(U_0^+,v_0^+),                                    \label{eq:Akh}\\
%   \F_{kh}(v)
%   :={}&\sum_{n=1}^N\int_{I_n}\ell(t;v)\dd t
%        +m(u_0,v_0^+).                                      \label{eq:Fkh}
% \end{align}
% The fully discrete primal solution is
% \begin{equation}
%   \boxed{\A_{kh}(U_{kh})(v_{kh})=\F_{kh}(v_{kh})
%   \quad\forall v_{kh}\in\X_{kh}^{r,p}.}
%   \label{eq:discrete-primal}
% \end{equation}



% \section{DWR Identity and Computable Enriched Weights}

% Extend \(\A_{kh}\) consistently to exact and enriched functions and define
% \begin{equation}
%   \rho_{kh}(U)(v):=\F_{kh}(v)-\A_{kh}(U)(v).
%   \label{eq:residual}
% \end{equation}
% The discrete adjoint is the backward-in-time problem
% \begin{equation}
%   \A_{kh}(w_{kh})(Z_{kh})=\J(w_{kh})
%   \quad\forall w_{kh}\in\X_{kh}^{r,p}.
%   \label{eq:discrete-adjoint}
% \end{equation}

% \begin{proposition}[Linear DWR identity]
% Assume that the primal problem, goal and discretisation are linear and consistent, and that quadrature and algebraic solves are exact. For every \(\varphi_{kh}\in\X_{kh}^{r,p}\),
% \begin{equation}
%   \boxed{\J(u)-\J(U_{kh})
%   =\rho_{kh}(U_{kh})(z)
%   =\rho_{kh}(U_{kh})(z-\varphi_{kh}).}
%   \label{eq:dwr-exact}
% \end{equation}
% \end{proposition}
% The first equality follows by testing the adjoint with \(u-U_{kh}\); the second uses Galerkin orthogonality \(\rho_{kh}(U_{kh})(\varphi_{kh})=0\). This is the linear DWR representation \citep{BeckerRannacher1996,BeckerRannacher2001}. For nonlinear operators or goals, the symmetric Lagrangian identity contains a dual residual, a factor \(1/2\), and a higher-order remainder \citep{BeckerRannacher2001}; those terms are absent from this linear specification.
% % 可以写进sec 2 in thesis

% Let \(Z^+\) be a numerical adjoint in a strictly enriched space \(\X_{kh}^+\supsetneq\X_{kh}^{r,p}\). The direct enriched-weight estimator is
% \begin{equation}
%   \eta_{\mathrm{joint}}
%   :=\rho_{kh}(U_{kh})(Z^+-Z_{kh}).
%   \label{eq:joint-estimator}
% \end{equation}
% This replaces the exact defect \(z-\varphi_{kh}\) by the numerical defect \(Z^+-Z_{kh}\). An enriched solve is more expensive but avoids relying solely on postprocessed interpolation. Reliability of enriched-space estimators is normally tied to a saturation assumption; interpolation recoveries need further recovery assumptions \citep{EndtmayerLangerWick2021}.
% % check



% \section{Separating Temporal and Spatial Goal Error}
% \label{sec:split}

% % \subsection{Split is prior to localisation}

% Let \(U_k\in\X_k^r(V)\) be the time-semidiscrete solution: time is discretised but space is still \(V\). Algebra gives
% \begin{equation}
%   \boxed{\J(u)-\J(U_{kh})
%   =\underbrace{\J(u)-\J(U_k)}_{e_k\ \text{(temporal)}}
%   +\underbrace{\J(U_k)-\J(U_{kh})}_{e_h\ \text{(spatial)}}.}
%   \label{eq:exact-split}
% \end{equation}
% % This splits the goal error between approximation levels; it does not split the operator into \(M\partial_t\) and \(\mathcal{L}\). 
% Separated temporal/spatial DWR representations based on this hierarchy were developed for parabolic problems and dynamic meshes by \citet{SchmichVexler2008}, and related formulations are used for nonstationary flow and modern space--time DWR methods \citep{BesierRannacher2012,ThieleWick2024,RothEtAl2024}.



% \newpage
\section{Automated Localisation by Bubble Recovery on Space--Time Cells}
\label{sec:localisation}

The bubble and cone functions are used to convert the weak residual into computable polynomial representations of its volume, spatial-facet, and temporal-interface parts. The DWR weight is the same enriched numerical dual defect:
\begin{equation}
 w:=Z^{+}-Z_{kh}.
 \label{eq:single-dual-weight}
\end{equation}
Here \(Z^{+}\) is one adjoint solution in a suitably enriched space and
\(Z_{kh}\) is the adjoint in the original discrete space. 

The localisation pipeline is therefore
\begin{equation}
 \boxed{
 \rho_{kh}(U_{kh})(\,\cdot\,)
 \xrightarrow{\text{bubble/cone projection}}
 \bigl(R^{\mathrm{vol}},R^{\mathrm{sf}},R^{\mathrm{tf}}\bigr)
 \xrightarrow{\text{single weight }w}
 \eta_{K,n}.}
 \label{eq:single-weight-bubble-pipeline}
\end{equation}
The indices \(\mathrm{sf}\) and \(\mathrm{tf}\) denote spatial-facet and temporal-facet residuals. Bubble recovery follows the stationary residual-representation construction of \citet{RognesLogg2013}.

For a general linear time-dependent problem, define
\begin{equation}
 \rho_{kh}(U_{kh})(V)
 :=
 \mathcal F_{kh}(V)-\mathcal A_{kh}(U_{kh},V).
 \label{eq:weak-residual-recovery}
\end{equation}
The global computable DWR estimator is
\begin{equation}
 \eta_{\mathrm{DWR}}^{\mathrm{glob}}
 :=
 \rho_{kh}(U_{kh})(w)
 =
 \rho_{kh}(U_{kh})(Z^{+}-Z_{kh}).
 \label{eq:single-weight-global-dwr}
\end{equation}
The purpose of the following construction is to express \cref{eq:single-weight-global-dwr} approximately as a signed sum over space--time cells \(Q_{K,n}=K\times I_n\).


\subsection{Stationary Bubble and Cone Projection}
\label{sec:stationary-residual-recovery}

Suppose that the local weak residual on a spatial cell \(K\) has a cell-and-facet representation
\begin{equation}
 \rho_K(v)
 =
 (R_K,v)_K
 +
 \sum_{F\subset\partial K}(R_{K,F},v)_F .
 \label{eq:stationary-residual-representation}
\end{equation}
The point of the method is to compute polynomial approximations to \(R_K\) and \(R_{K,F}\) from the weak residual without deriving their strong formulas by hand.

Let \(K\) be a \(d\)-simplex with barycentric coordinates \(\lambda_1^K,\ldots,\lambda_{d+1}^K\), and define the cell bubble
\begin{equation}
 b_K:=\prod_{i=1}^{d+1}\lambda_i^K,
 \qquad b_K|_{\partial K}=0.
 \label{eq:recovery-cell-bubble}
\end{equation}
For a polynomial recovery space \(Q_K^p\), find \(R_K^B\in Q_K^p\) such that
\begin{equation}
 \boxed{
 (R_K^B,b_K\phi)_K
 =
 \rho_K(b_K\phi)
 \qquad\forall\phi\in Q_K^p.}
 \label{eq:stationary-cell-projection}
\end{equation}
Because \(b_K\phi\) is zero on every facet, substitution of \cref{eq:stationary-residual-representation} gives
\[
 \rho_K(b_K\phi)=(R_K,b_K\phi)_K.
\]

To recover a facet residual, let \(F\subset\partial K\) and introduce a cone
\begin{equation}
 \beta_{K,F}:=\prod_{i\in\mathcal I_F}\lambda_i^K,
 \label{eq:recovery-facet-cone}
\end{equation}
where \(\mathcal I_F\) contains the vertices belonging to \(F\). The cone is nonzero on \(F\) but vanishes on every other facet. Therefore
\begin{equation}
 \rho_K(\beta_{K,F}\psi)
 =
 (R_K,\beta_{K,F}\psi)_K
 +
 (R_{K,F},\beta_{K,F}\psi)_F.
 \label{eq:facet-probe-expanded}
\end{equation}
Hence find
\(R_{K,F}^B\in Q_F^q\) from
\begin{equation}
 \boxed{
 (R_{K,F}^B,\beta_{K,F}\psi)_F
 =
 \rho_K(\beta_{K,F}\psi)
 -
 (R_K^B,\beta_{K,F}\psi)_K
 \qquad\forall\psi\in Q_F^q.}
 \label{eq:stationary-facet-projection}
\end{equation}


\subsection{Extension: Treating Time as An Additional Mesh Dimension}
\label{sec:space-time-residual-recovery}

Let
\[
 I_n=(t_{n-1},t_n],
 \qquad
 k_n=t_n-t_{n-1},
 \qquad
 Q_{K,n}:=K\times I_n,
 \qquad
 \tau:=\frac{t-t_{n-1}}{k_n}.
\]
The space--time cell \(Q_{K,n}\) has three kinds of residual entities:
\begin{equation}
 \rho_{K,n}(V)
 =
 (R_{K,n}^{\mathrm{vol}},V)_{Q_{K,n}}
 +
 \sum_{F\subset\partial K}
 (R_{K,F,n}^{\mathrm{sf}},V)_{F\times I_n}
 +
 (R_{K,n-1}^{\mathrm{tf}},V(t_{n-1}^{+}))_K .
 \label{eq:space-time-residual-representation}
\end{equation}
For a dG-in-time method, the temporal-facet term contains the state jump at \(t_{n-1}\); for \(n=1\), it contains the initial-condition mismatch. On changing spatial meshes it must also contain the transfer defect \cite{ErikssonJohnsonThomee1985,SchmichVexler2008}.

\paragraph{Space--time volume projection.}
Consider the temporal interior bubble function
\begin{equation}
 b_n^t(\tau):=4\tau(1-\tau).
 \label{eq:temporal-interior-bubble}
\end{equation}

Define the space--time interior bubble
\[
 B_{K,n}(x,t):=b_K(x)b_n^t(t).
\]
It vanishes on every spatial and temporal face of \(Q_{K,n}\). The recovered volume density \(R_{K,n}^{B,\mathrm{vol}}\) is defined by
\begin{equation}
 \boxed{
 (R_{K,n}^{B,\mathrm{vol}},B_{K,n}\Phi)_{Q_{K,n}}
 =
 \rho_{K,n}(B_{K,n}\Phi)
 \qquad
 \forall\Phi\in Q_{K,n}^{p_x,p_t}.}
 \label{eq:space-time-volume-projection}
\end{equation}


\paragraph{Spatial-facet projection.}
For \(F\subset\partial K\), use
\begin{equation}
 C_{K,F,n}^{\mathrm{sf}}(x,t)
 :=
 \beta_{K,F}(x)b_n^t(t).
 \label{eq:space-time-spatial-cone}
\end{equation}
This function is nonzero only on the selected spatial face \(F\times I_n\) and vanishes at both time endpoints. After the volume solve, recover \(R_{K,F,n}^{B,\mathrm{sf}}\) from
\begin{align}
 &(R_{K,F,n}^{B,\mathrm{sf}},
       C_{K,F,n}^{\mathrm{sf}}\Psi)_{F\times I_n}
 \notag\\
 &\qquad =
 \rho_{K,n}(C_{K,F,n}^{\mathrm{sf}}\Psi)
 -
 (R_{K,n}^{B,\mathrm{vol}},
       C_{K,F,n}^{\mathrm{sf}}\Psi)_{Q_{K,n}}
 \quad\forall\Psi\in Q_{F,n}^{q_x,q_t}.
 \label{eq:space-time-spatial-facet-projection}
\end{align}

\paragraph{Temporal-facet projection.}
% The interior bubble \(b_n^t\) cannot see a temporal jump because it is zero at \(t_{n-1}\) and \(t_n\).  To probe the lower temporal face, 
We use the temporal cone
\begin{equation}
 C_{K,n}^{\mathrm{tf},-}(x,t)
 :=
 b_K(x)(1-\tau).
 \label{eq:temporal-trace-cone}
\end{equation}
It is nonzero at \(t_{n-1}^{+}\), zero at \(t_n^{-}\), and zero on every spatial facet. Consequently,
\begin{align}
 &(R_{K,n-1}^{B,\mathrm{tf}},b_K\chi)_K
 \notag\\
 &\qquad =
 \rho_{K,n}(C_{K,n}^{\mathrm{tf},-}\chi)
 -
 (R_{K,n}^{B,\mathrm{vol}},
       C_{K,n}^{\mathrm{tf},-}\chi)_{Q_{K,n}}
 \qquad\forall\chi\in Q_K^{q_t}.
 \label{eq:temporal-facet-projection}
\end{align}


\subsection{Signed Indicator for Every Space--Time Cell}
\label{sec:space-time-cell-indicators}

The original numerical dual defect \(w=Z^+-Z_{kh}\) weights every recovered residual entity:
\begin{align}
 \eta_{K,n}^{\mathrm{vol}}
 &:=
 (R_{K,n}^{B,\mathrm{vol}},w)_{Q_{K,n}},
 \label{eq:single-volume-contribution}\\
 \eta_{F,n}^{\mathrm{sf}}
 &:=
 (R_{F,n}^{B,\mathrm{sf}},w)_{F\times I_n},
 \label{eq:single-spatial-facet-contribution}\\
 \eta_{K,n-1}^{\mathrm{tf}}
 &:=
 (R_{K,n-1}^{B,\mathrm{tf}},w(t_{n-1}^{+}))_K .
 \label{eq:single-temporal-facet-contribution}
\end{align}

% An interior spatial facet is one residual entity shared by two neighbouring space--time cells. Meanwhile, time is used as an additional coordinate for constructing and marking space--time cells. The jump at \(t_{n-1}\) is tested with the right trace \(w(t_{n-1}^{+})\) and its entire signed contribution is assigned to the following cell slab \(Q_{K,n}=K\times I_n\).

Let \(\mathcal E_{K,n}^{\mathrm{sf,int}}\) denote the interior spatial
facets of \(Q_{K,n}\). The signed local indicator is
\begin{equation}
 \boxed{
 \eta_{K,n}
 :=
 \eta_{K,n}^{\mathrm{vol}}
 +
 \frac12\sum_{F\in\mathcal E_{K,n}^{\mathrm{sf,int}}}
       \eta_{F,n}^{\mathrm{sf}}
 +
 \sum_{F\subset\partial K\cap\partial\Omega}
       \eta_{F,n}^{\mathrm{bdry}}
 +
 \eta_{K,n-1}^{\mathrm{tf}}.}
 \label{eq:single-space-time-cell-indicator}
\end{equation}
Every unique spatial interior facet is counted once through its two half-contributions, and every temporal interface is counted once through its full contribution to the following slab.

\begin{align}
 \eta_n^{\mathrm{slab}}
 &:=
 \sum_{K\in\mathcal T_h^n}\eta_{K,n},
 \label{eq:signed-slab-indicator}\\
 s_n^{\mathrm{slab}}
 &:=
 \sum_{K\in\mathcal T_h^n}|\eta_{K,n}|.
 \label{eq:absolute-slab-score}
\end{align}
The signed quantity \(\eta_n^{\mathrm{slab}}\) is useful for checking how each slab contributes to the global goal estimate. The nonnegative quantity \(s_n^{\mathrm{slab}}\) measures the total marking activity in \(I_n\) without cancellation.

The recovered signed global estimator is
\begin{equation}
 \widehat\eta_B
 :=
 \sum_{n=1}^N\sum_{K\in\mathcal T_h^n}\eta_{K,n}
 =
 \sum_{n=1}^N\eta_n^{\mathrm{slab}}.
 \label{eq:single-recovered-global-estimator}
\end{equation}
It should be compared against the directly assembled value \(\eta_{\mathrm{DWR}}^{\mathrm{glob}}\) in \cref{eq:single-weight-global-dwr}. Their discrepancy measures the localisation error introduced by the finite-degree residual projections.


\subsection{Marking and Refinement}
\label{sec:joint-space-time-marking}

Use
\begin{equation}
 s_{K,n}:=|\eta_{K,n}|
 \label{eq:single-cell-marking-score}
\end{equation}
as the nonnegative marking score. Select an approximately minimal set
\(\mathcal M_{hk}\) satisfying
\begin{equation}
 \sum_{(K,n)\in\mathcal M_{hk}}s_{K,n}
 \ge
 \theta\sum_{n=1}^N\sum_{K\in\mathcal T_h^n}s_{K,n},
 \qquad 0<\theta\le1.
 \label{eq:space-time-dorfler}
\end{equation}
This is D\"orfler bulk marking applied directly to space--time cell contributions \cite{Doerfler1996}.

Each marked pair \((K,n)\in\mathcal M_{hk}\) triggers both actions:
\begin{enumerate}[leftmargin=*]
 \item mark \(K\) for spatial refinement on the mesh associated with \(I_n\);
 \item mark \(I_n\) for temporal bisection.
\end{enumerate}
Equivalently, construct
\begin{align}
 \mathcal M_h^n
 &:=
 \{K\in\mathcal T_h^n:(K,n)\in\mathcal M_{hk}\},
 \label{eq:spatial-mark-projection}\\
 \mathcal M_t
 &:=
 \{n:\mathcal M_h^n\neq\varnothing\}.
 \label{eq:temporal-mark-projection}
\end{align}
Every \(I_n\) appearing in \(\mathcal M_t\) is bisected exactly once. The spatial marked set is the union of all marked \(K\)'s in that slab. After the bisection, the selected spatial refinement is inherited by both child intervals.


\subsection{Automated Adaptive Loop}
\label{sec:single-indicator-loop}

The resulting solve--estimate--mark--refine cycle is:
\begin{enumerate}[leftmargin=*]
 \item Solve the primal problem forward for \(U_{kh}\).
 \item Solve the original discrete adjoint for \(Z_{kh}\) and one enriched adjoint for \(Z^+\); form \(w=Z^+-Z_{kh}\).
 \item Pass the weak residual \(\rho_{kh}(U_{kh})(\,\cdot\,)\) to the recovery module.
 \item Recover the volume residual, then the spatial-facet residuals, and the temporal-interface residuals.
 \item Weight every recovered entity by \(w\), assemble the signed \(\eta_{K,n}\), and check \cref{eq:single-recovered-global-estimator} against \cref{eq:single-weight-global-dwr}.
 \item Stop if \(|\eta_{\mathrm{DWR}}^{\mathrm{glob}}|\le\mathrm{TOL}_J\).
 \item Otherwise mark space--time cells by \cref{eq:space-time-dorfler}; bisect every selected time interval once and refine the selected spatial cells on both resulting child slabs.
 \item Rebuild meshes, transfer operators and repeat.
\end{enumerate}



\newpage
\bibliographystyle{unsrtnat}
\bibliography{automated_dwr_time_adaptivity}

\end{document}
